% Pré-preenchimento: revisar antes da entrega.
\section{Ameaças à Validade}
\label{sec:ameacas}

\textcolor{orange}{\textbf{Validade de construto.} Não há, neste estudo, um padrão-ouro humano de qualidade das descrições: a precisão da avaliação é inferida a partir do mecanismo pelo qual a nota muda (RQ1), não medida contra um julgamento de referência. Assim, demonstra-se que a mudança de nota tem causa identificável no código, mas não que a nota resultante seja, em sentido absoluto, mais correta. Além disso, tanto as avaliações quanto as etapas da análise qualitativa (síntese das justificativas, indução de categorias e classificação) são conduzidas por LLMs, e a verificação manual compara a categoria atribuída com a justificativa sintetizada e o \textit{reasoning} dos juízes, não diretamente com o código-fonte. O Gemini 3.5 Flash-Lite atua como juiz e também como sintetizador das justificativas, o que pode favorecer a sua própria leitura dos casos.}

\textcolor{orange}{\textbf{Validade interna.} O código fornecido ao avaliador é obtido por uma heurística de resolução de chamadas baseada em nome, não em tipos, limitada a três níveis de profundidade e a 32.000 caracteres (Apêndice~\ref{apendice:secao}). Parte da lógica relevante pode, portanto, ficar de fora, ou trechos de funções homônimas podem ser incluídos por engano, embora resoluções ambíguas sejam sinalizadas. As respostas dos juízes não são determinísticas, e a classificação das justificativas em categorias foi confirmada em 72\% da amostra verificada, com discordâncias concentradas entre categorias de fronteira próxima. Por outro lado, a amostra de divergência confirmada exige que todos os juízes concordem que a nota mudou, o que a torna conservadora frente ao ruído de um juiz isolado.}

\textcolor{orange}{\textbf{Validade externa.} A amostra restringe-se a 206 repositórios públicos do GitHub com ao menos 100 estrelas, concentrados em Python e TypeScript, e três repositórios respondem por 46,6\% das ferramentas, o que pode enviesar resultados agregados. Os resultados também refletem três modelos específicos em uma versão e momento determinados; como os próprios juízes divergem na direção do efeito (RQ2), não é possível generalizar o comportamento observado para outros avaliadores.}

\textcolor{orange}{\textbf{Validade de conclusão.} Com milhares de pares por atributo, até diferenças pequenas resultam em $p$-valores muito baixos; por isso, o Teste de Wilcoxon é acompanhado da transição de quartil, que dimensiona o efeito. Não foi aplicada correção para comparações múltiplas, mas, como todos os $p$-valores são inferiores a $10^{-80}$, uma correção de Bonferroni para seis testes ($\alpha = 0{,}0083$) não alteraria nenhuma conclusão. Por fim, a média entre juízes, usada na RQ3, pode atenuar efeitos de sentidos opostos, como a redução de notas pelo DeepSeek frente ao aumento pelo Gemini e pelo Qwen3, razão pela qual a RQ2 apresenta os resultados por juiz.}
